\RequirePackage{fix-cm}
\documentclass[12pt,letterpaper]{article}
\usepackage[margin=.5in]{geometry}
\usepackage{tikz}
\usetikzlibrary{arrows.meta,patterns}
\pdfmapfile{+cm.map}
\pagestyle{empty}
\setlength{\parindent}{0pt}
\begin{document}
\null\begin{tikzpicture}[remember picture,overlay]
\begin{scope}[shift={(current page.north west)},x=1cm,y=-1cm]
\begin{scope}[shift={(1.5,1.5)}]
\node[anchor=north west,inner sep=0,align=left,text width=18.5cm,font=\sffamily\fontsize{11}{13.75}\selectfont] at (0,-0.5) {Week 35 / Footprint borders / Grades 2--3};
\node[anchor=north west,inner sep=0,align=left,text width=17cm,font=\sffamily\fontsize{10}{12.5}\selectfont] at (0,24.7) {Bellingham Math Circle / Week 35 / W35-grades-2-3-v2};
\node[anchor=north west,inner sep=0,align=left,text width=0.5cm,font=\sffamily\fontsize{10}{12.5}\selectfont] at (18,24.7) {1};
\node[anchor=north west,inner sep=0,align=left,text width=18.4cm,font=\sffamily\fontsize{14}{17.5}\selectfont] at (0,0.3) {Use the loose footprints and long strips. Each border repeats forever both ways. A match must fit every footprint. Only footprints count, not strip ends or guide lines. A slide must move the copy. Test with a tracing copy; leave the original in place.};
\draw[dashed,gray] (0.5,4.9) -- (4.7,4.9);
\draw[dashed,gray] (7.1,4.9) -- (11.299999999999999,4.9);
\draw[dashed,gray] (13.700000000000001,4.9) -- (17.900000000000002,4.9);
\begin{scope}[shift={(2.6,3.6)},rotate=0,xscale=0.68,yscale=0.68]
\draw[line width=.7pt,fill=gray!15] (-1.3,-1.1) -- (1.3,-1.1) -- (1.3,-.4) -- (.8,-.4) -- (.8,-.65) -- (-.5,-.65) -- (-.5,0) -- (.6,0) -- (.6,.45) -- (-.5,.45) -- (-.5,1.1) -- (-1.3,1.1) -- cycle;
\draw[-{Stealth},line width=.65pt] (-.25,-.86) -- (.55,-.86);\end{scope}
\begin{scope}[shift={(9.2,6.2)},rotate=0,xscale=0.68,yscale=-0.68]
\draw[line width=.7pt,fill=gray!15] (-1.3,-1.1) -- (1.3,-1.1) -- (1.3,-.4) -- (.8,-.4) -- (.8,-.65) -- (-.5,-.65) -- (-.5,0) -- (.6,0) -- (.6,.45) -- (-.5,.45) -- (-.5,1.1) -- (-1.3,1.1) -- cycle;
\draw[-{Stealth},line width=.65pt] (-.25,-.86) -- (.55,-.86);\end{scope}
\begin{scope}[shift={(16.8,6.2)},rotate=0,xscale=0.68,yscale=-0.68]
\draw[line width=.7pt,fill=gray!15] (-1.3,-1.1) -- (1.3,-1.1) -- (1.3,-.4) -- (.8,-.4) -- (.8,-.65) -- (-.5,-.65) -- (-.5,0) -- (.6,0) -- (.6,.45) -- (-.5,.45) -- (-.5,1.1) -- (-1.3,1.1) -- cycle;
\draw[-{Stealth},line width=.65pt] (-.25,-.86) -- (.55,-.86);\end{scope}
\node[anchor=north west,inner sep=0,align=left,text width=0.7cm,font=\sffamily\fontsize{12}{15.0}\selectfont] at (1.05,3.3000000000000003) {A};
\node[anchor=north west,inner sep=0,align=left,text width=0.7cm,font=\sffamily\fontsize{12}{15.0}\selectfont] at (7.6499999999999995,5.9) {A};
\node[anchor=north west,inner sep=0,align=left,text width=0.7cm,font=\sffamily\fontsize{12}{15.0}\selectfont] at (15.25,5.9) {A};
\draw[-{Stealth},line width=.9pt] (4.8,4.9) -- (6.8,4.9);
\draw[-{Stealth},line width=.9pt] (11.5,4.9) -- (13.4,4.9);
\node[anchor=north west,inner sep=0,align=left,text width=4.5cm,font=\sffamily\fontsize{12}{15.0}\selectfont] at (0.7,7.4) {start};
\node[anchor=north west,inner sep=0,align=left,text width=5cm,font=\sffamily\fontsize{12}{15.0}\selectfont] at (7.1,7.4) {flip over the line};
\node[anchor=north west,inner sep=0,align=left,text width=4.7cm,font=\sffamily\fontsize{12}{15.0}\selectfont] at (13.7,7.4) {then slide right};
\draw[-{Stealth},line width=.8pt] (14.7,7.1) -- (15.7,7.1);
\node[anchor=north west,inner sep=0,align=left,text width=18.4cm,font=\sffamily\fontsize{15}{18.75}\selectfont] at (0,8.9) {\textbf{Problem 1:} Make two repeating borders whose smallest matching slides are different lengths. Mark one matching slide on each tracing copy.};
\draw[gray!50,line width=.6pt] (.35,11.2) rectangle (18,17.2);
\draw[dashed,line width=.6pt] (.35,14.2) -- (18,14.2);
\draw[-{Stealth},line width=.9pt] (0.6,14.2) -- (0,14.2);
\draw[-{Stealth},line width=.9pt] (17.799999999999997,14.2) -- (18.4,14.2);
\draw[gray!45] (1.7,14.08) -- (1.7,14.319999999999999);
\draw[gray!45] (4.7,14.08) -- (4.7,14.319999999999999);
\draw[gray!45] (7.7,14.08) -- (7.7,14.319999999999999);
\draw[gray!45] (10.7,14.08) -- (10.7,14.319999999999999);
\draw[gray!45] (13.7,14.08) -- (13.7,14.319999999999999);
\draw[gray!45] (16.7,14.08) -- (16.7,14.319999999999999);
\draw[gray!50,line width=.6pt] (.35,18.1) rectangle (18,24.1);
\draw[dashed,line width=.6pt] (.35,21.1) -- (18,21.1);
\draw[-{Stealth},line width=.9pt] (0.6,21.1) -- (0,21.1);
\draw[-{Stealth},line width=.9pt] (17.799999999999997,21.1) -- (18.4,21.1);
\draw[gray!45] (1.7,20.98) -- (1.7,21.220000000000002);
\draw[gray!45] (4.7,20.98) -- (4.7,21.220000000000002);
\draw[gray!45] (7.7,20.98) -- (7.7,21.220000000000002);
\draw[gray!45] (10.7,20.98) -- (10.7,21.220000000000002);
\draw[gray!45] (13.7,20.98) -- (13.7,21.220000000000002);
\draw[gray!45] (16.7,20.98) -- (16.7,21.220000000000002);\end{scope}\end{scope}\end{tikzpicture}
\newpage
\null\begin{tikzpicture}[remember picture,overlay]
\begin{scope}[shift={(current page.north west)},x=1cm,y=-1cm]
\begin{scope}[shift={(1.5,1.5)}]
\node[anchor=north west,inner sep=0,align=left,text width=18.5cm,font=\sffamily\fontsize{11}{13.75}\selectfont] at (0,-0.5) {Week 35 / Footprint borders / Grades 2--3};
\node[anchor=north west,inner sep=0,align=left,text width=17cm,font=\sffamily\fontsize{10}{12.5}\selectfont] at (0,24.7) {Bellingham Math Circle / Week 35 / W35-grades-2-3-v2};
\node[anchor=north west,inner sep=0,align=left,text width=0.5cm,font=\sffamily\fontsize{10}{12.5}\selectfont] at (18,24.7) {2};
\node[anchor=north west,inner sep=0,align=left,text width=18.4cm,font=\sffamily\fontsize{15}{18.75}\selectfont] at (0,0.5) {\textbf{Problem 2:} A glide is a flip over the middle line followed by a nonzero slide along it. Make a border that matches after a glide, but fails after the flip alone.};
\draw[gray!50,line width=.6pt] (.35,4.1) rectangle (18,10.1);
\draw[dashed,line width=.6pt] (.35,7.1) -- (18,7.1);
\draw[-{Stealth},line width=.9pt] (0.6,7.1) -- (0,7.1);
\draw[-{Stealth},line width=.9pt] (17.799999999999997,7.1) -- (18.4,7.1);
\draw[gray!45] (1.7,6.9799999999999995) -- (1.7,7.22);
\draw[gray!45] (4.7,6.9799999999999995) -- (4.7,7.22);
\draw[gray!45] (7.7,6.9799999999999995) -- (7.7,7.22);
\draw[gray!45] (10.7,6.9799999999999995) -- (10.7,7.22);
\draw[gray!45] (13.7,6.9799999999999995) -- (13.7,7.22);
\draw[gray!45] (16.7,6.9799999999999995) -- (16.7,7.22);
\draw[gray!50,line width=.6pt] (.35,12.0) rectangle (18,18.0);
\draw[dashed,line width=.6pt] (.35,15) -- (18,15);
\draw[-{Stealth},line width=.9pt] (0.6,15) -- (0,15);
\draw[-{Stealth},line width=.9pt] (17.799999999999997,15) -- (18.4,15);
\draw[gray!45] (1.7,14.88) -- (1.7,15.12);
\draw[gray!45] (4.7,14.88) -- (4.7,15.12);
\draw[gray!45] (7.7,14.88) -- (7.7,15.12);
\draw[gray!45] (10.7,14.88) -- (10.7,15.12);
\draw[gray!45] (13.7,14.88) -- (13.7,15.12);
\draw[gray!45] (16.7,14.88) -- (16.7,15.12);
\node[anchor=north west,inner sep=0,align=left,text width=18.4cm,font=\sffamily\fontsize{15}{18.75}\selectfont] at (0,20.3) {\textbf{Problem 3:} On your border, how does the shortest matching glide compare with the shortest matching slide?};\end{scope}\end{scope}\end{tikzpicture}
\newpage
\null\begin{tikzpicture}[remember picture,overlay]
\begin{scope}[shift={(current page.north west)},x=1cm,y=-1cm]
\begin{scope}[shift={(1.5,1.5)}]
\node[anchor=north west,inner sep=0,align=left,text width=18.5cm,font=\sffamily\fontsize{11}{13.75}\selectfont] at (0,-0.5) {Week 35 / Footprint borders / Grades 2--3};
\node[anchor=north west,inner sep=0,align=left,text width=17cm,font=\sffamily\fontsize{10}{12.5}\selectfont] at (0,24.7) {Bellingham Math Circle / Week 35 / W35-grades-2-3-v2};
\node[anchor=north west,inner sep=0,align=left,text width=0.5cm,font=\sffamily\fontsize{10}{12.5}\selectfont] at (18,24.7) {3};
\node[anchor=north west,inner sep=0,align=left,text width=18.4cm,font=\sffamily\fontsize{14}{17.5}\selectfont] at (0,0.3) {A half-turn turns the copy halfway around a chosen point on the table. You may choose any point.};
\node[anchor=north west,inner sep=0,align=left,text width=18.4cm,font=\sffamily\fontsize{15}{18.75}\selectfont] at (0,2.1) {\textbf{Problem 4:} Make a repeating border that matches after a flip over the middle line. Can you keep it from matching after any half-turn?};
\draw[gray!50,line width=.6pt] (.35,5.300000000000001) rectangle (18,11.3);
\draw[dashed,line width=.6pt] (.35,8.3) -- (18,8.3);
\draw[-{Stealth},line width=.9pt] (0.6,8.3) -- (0,8.3);
\draw[-{Stealth},line width=.9pt] (17.799999999999997,8.3) -- (18.4,8.3);
\draw[gray!45] (1.7,8.180000000000001) -- (1.7,8.42);
\draw[gray!45] (4.7,8.180000000000001) -- (4.7,8.42);
\draw[gray!45] (7.7,8.180000000000001) -- (7.7,8.42);
\draw[gray!45] (10.7,8.180000000000001) -- (10.7,8.42);
\draw[gray!45] (13.7,8.180000000000001) -- (13.7,8.42);
\draw[gray!45] (16.7,8.180000000000001) -- (16.7,8.42);
\draw[gray!50,line width=.6pt] (.35,13.7) rectangle (18,19.7);
\draw[dashed,line width=.6pt] (.35,16.7) -- (18,16.7);
\draw[-{Stealth},line width=.9pt] (0.6,16.7) -- (0,16.7);
\draw[-{Stealth},line width=.9pt] (17.799999999999997,16.7) -- (18.4,16.7);
\draw[gray!45] (1.7,16.58) -- (1.7,16.82);
\draw[gray!45] (4.7,16.58) -- (4.7,16.82);
\draw[gray!45] (7.7,16.58) -- (7.7,16.82);
\draw[gray!45] (10.7,16.58) -- (10.7,16.82);
\draw[gray!45] (13.7,16.58) -- (13.7,16.82);
\draw[gray!45] (16.7,16.58) -- (16.7,16.82);\end{scope}\end{scope}\end{tikzpicture}
\newpage
\null\begin{tikzpicture}[remember picture,overlay]
\begin{scope}[shift={(current page.north west)},x=1cm,y=-1cm]
\begin{scope}[shift={(1.5,1.5)}]
\node[anchor=north west,inner sep=0,align=left,text width=18.5cm,font=\sffamily\fontsize{11}{13.75}\selectfont] at (0,-0.5) {Week 35 / Footprint borders / Grades 2--3};
\node[anchor=north west,inner sep=0,align=left,text width=17cm,font=\sffamily\fontsize{10}{12.5}\selectfont] at (0,24.7) {Bellingham Math Circle / Week 35 / W35-grades-2-3-v2};
\node[anchor=north west,inner sep=0,align=left,text width=0.5cm,font=\sffamily\fontsize{10}{12.5}\selectfont] at (18,24.7) {4};
\node[anchor=north west,inner sep=0,align=left,text width=18.4cm,font=\sffamily\fontsize{15}{18.75}\selectfont] at (0,0.5) {\textbf{Problem 5:} Make a repeating border that matches after a half-turn, but fails after the flip over the middle line.};
\draw[gray!50,line width=.6pt] (.35,4.1) rectangle (18,10.1);
\draw[dashed,line width=.6pt] (.35,7.1) -- (18,7.1);
\draw[-{Stealth},line width=.9pt] (0.6,7.1) -- (0,7.1);
\draw[-{Stealth},line width=.9pt] (17.799999999999997,7.1) -- (18.4,7.1);
\draw[gray!45] (1.7,6.9799999999999995) -- (1.7,7.22);
\draw[gray!45] (4.7,6.9799999999999995) -- (4.7,7.22);
\draw[gray!45] (7.7,6.9799999999999995) -- (7.7,7.22);
\draw[gray!45] (10.7,6.9799999999999995) -- (10.7,7.22);
\draw[gray!45] (13.7,6.9799999999999995) -- (13.7,7.22);
\draw[gray!45] (16.7,6.9799999999999995) -- (16.7,7.22);
\draw[gray!50,line width=.6pt] (.35,13.100000000000001) rectangle (18,19.1);
\draw[dashed,line width=.6pt] (.35,16.1) -- (18,16.1);
\draw[-{Stealth},line width=.9pt] (0.6,16.1) -- (0,16.1);
\draw[-{Stealth},line width=.9pt] (17.799999999999997,16.1) -- (18.4,16.1);
\draw[gray!45] (1.7,15.980000000000002) -- (1.7,16.220000000000002);
\draw[gray!45] (4.7,15.980000000000002) -- (4.7,16.220000000000002);
\draw[gray!45] (7.7,15.980000000000002) -- (7.7,16.220000000000002);
\draw[gray!45] (10.7,15.980000000000002) -- (10.7,16.220000000000002);
\draw[gray!45] (13.7,15.980000000000002) -- (13.7,16.220000000000002);
\draw[gray!45] (16.7,15.980000000000002) -- (16.7,16.220000000000002);
\node[anchor=north west,inner sep=0,align=left,text width=18.4cm,font=\sffamily\fontsize{15}{18.75}\selectfont] at (0,21.2) {\textbf{Problem 6:} Move footprints in one of your borders so that an old matching move fails but the new border still repeats. You may use a larger repeating part.};\end{scope}\end{scope}\end{tikzpicture}
\end{document}
